\section{Introduction}

Autoregressive decoding in large language models (LLMs) is inherently
sequential: each generated token requires a full forward pass of the target
model. This constraint directly limits single-stream throughput in interactive
inference settings. Speculative decoding mitigates this bottleneck by adding a
smaller draft model that proposes a length-$k$ token block, followed by one
target-side verification step; modified rejection sampling preserves the target
distribution~\cite{Leviathan2023, Chen2023specsampling}.

Although speculative decoding is now well studied, deployment guidance remains
uncertain for consumer-GPU settings. Reported gains in the literature are often
obtained on data-center hardware and are sensitive to model pairing, proposal
length, and sampling regime. For practitioners operating on a single RTX-class
GPU, the key questions are practical: which draft size is worthwhile, which
$k$ is optimal, and whether deterministic versus stochastic decoding changes
the speed--quality frontier.

This study addresses three practical gaps: cross-model and cross-hardware
validation, sample- and seed-level variability reporting, and statistical
testing against matched autoregressive baselines. The goal is a reproducible,
hardware-aware characterization rather than a new decoding algorithm.

\subsection{Research questions}
\label{sec:rqs}

We address four research questions:
\begin{enumerate}\itemsep0pt
  \item[\textbf{RQ1.}] How much end-to-end speedup does vanilla
        speculative decoding deliver against autoregressive decoding on
        consumer RTX hardware with Qwen2.5 and Qwen3 pairings?
  \item[\textbf{RQ2.}] How do draft length and acceptance jointly determine
        net latency on consumer GPUs?
  \item[\textbf{RQ3.}] How sensitive is the speed--quality trade-off to
        deterministic versus stochastic decoding?
  \item[\textbf{RQ4.}] Do policy rankings and absolute gains transfer from
        desktop RTX4090 to RTX5090-laptop?
\end{enumerate}

\subsection{Contributions}
\label{sec:contributions}

This paper makes the following contributions:
\begin{enumerate}\itemsep0pt
  \item A reproducible protocol using frozen manifests and matched prompts on
        GSM8K, MMLU, and CNN/DailyMail.
  \item A fixed-policy study of Qwen2.5 and Qwen3 with
        $k\in\{4,8,16\}$ under deterministic and stochastic decoding.
  \item Cross-hardware validation showing that policy ranking transfers from
        RTX4090 to RTX5090-laptop while absolute speedup does not.
  \item Sample-level dispersion, targeted repeat-seed runs, paired tests with
        Holm--Bonferroni correction, and effect sizes against matched baselines.
\end{enumerate}

Section~\ref{sec:related} reviews prior work; Sections~\ref{sec:protocol}--
\ref{sec:impl} define the protocol, metrics, and implementation; and
Sections~\ref{sec:results}--\ref{sec:discussion} present results, limitations,
and conclusions.
